\documentclass[10pt,conference,compsocconf]{IEEEtran}

% ============================================================
% PACKAGES
% ============================================================
\usepackage{amsmath,amssymb,amsfonts}
\usepackage{graphicx}
\usepackage{booktabs}
\usepackage{microtype}
\usepackage{hyperref}
\usepackage{array}
\usepackage{balance}

% ============================================================
% COMPACT IEEE LAYOUT
% ============================================================
\setlength{\textfloatsep}{5pt plus 1pt minus 2pt}
\setlength{\floatsep}{5pt plus 1pt minus 2pt}
\setlength{\intextsep}{5pt plus 1pt minus 2pt}
\setlength{\abovecaptionskip}{2pt}
\setlength{\belowcaptionskip}{0pt}

\begin{document}

% ============================================================
% TITLE
% ============================================================
\title{Predicting Cardiovascular Disease Risk from Survey Health Indicators:
A Regularized Logistic and Feature-Engineering Approach}

\author{
\IEEEauthorblockN{Team CS-433}
\IEEEauthorblockA{
\textit{School of Computer and Communication Sciences, EPFL, Switzerland}
}
}

\maketitle

% ============================================================
% ABSTRACT
% ============================================================
\begin{abstract}
We investigate regularized logistic regression for predicting myocardial
infarction and coronary heart disease from the 2015 CDC Behavioral Risk Factor
Surveillance System (BRFSS). The course dataset contains 328,135 training
respondents and 321 input features. Modeling is implemented in NumPy and the
Python standard library. Candidate representations include missingness
indicators, quantile bins, rank transforms, piecewise-linear splines, and
clinical interactions. The logistic-only search reached a validation $F_1$ of
$0.449939$ on a stratified training-data holdout. This is an interim tuning
score: the search stopped before final cross-validation was completed, and
final ensemble results are not yet available.
\end{abstract}

% ============================================================
% INTRODUCTION
% ============================================================
\section{Introduction}

The CDC Behavioral Risk Factor Surveillance System (BRFSS)~\cite{cdc2015} collects
self-reported information on chronic conditions, health behaviors, and
preventive practices. Our objective is to predict whether a respondent has
coronary heart disease or myocardial infarction,
$\text{MICHD}\in\{-1,+1\}$.

The task presents four main difficulties. First, the positive prevalence is
only approximately $8.83\%$, so a trivial negative classifier already obtains
about $91.17\%$ accuracy while having $F_1=0$; we therefore use $F_1$ as the
primary selection metric. Second, BRFSS contains administrative response codes:
for example, \texttt{88} may represent ``none'', whereas \texttt{77} and
\texttt{99} denote unavailable responses. Third, all learning algorithms must
be implemented using only NumPy and the Python standard library. Finally, all
imputation, scaling, spline knots, and decision thresholds must be estimated
from fitting rows, with decision thresholds selected using validation labels.

% ============================================================
% PREPROCESSING
% ============================================================
\section{Preprocessing \& Feature Engineering}

\subsection{CDC Decoding and Cleaning}

Survey identifiers and date fields such as \texttt{SEQNO}, \texttt{\_PSU},
and \texttt{IDATE} are discarded because they carry no meaningful clinical
signal. For \texttt{PHYSHLTH}, \texttt{MENTHLTH}, and \texttt{POORHLTH},
which measure unhealthy days during the previous month, we map
$88\rightarrow0$ and convert \texttt{77}/\texttt{99} to missing values.

BRFSS questionnaire logic is also used to distinguish missingness from
administrative skips. In particular, when both physical and mental unhealthy
days are zero, an omitted \texttt{POORHLTH} response corresponds to a valid
skip and is therefore imputed as zero. Other missing numerical values are
imputed using training-fold medians, with an additional binary missingness
indicator for each affected feature.

\subsection{Feature Representations}

We compare several representations. Continuous variables are clipped to their
training $[0.5\%,99.5\%]$ quantiles and standardized. \textbf{Quantile bins}
capture threshold effects, while \textbf{empirical rank scaling} maps
variables to their empirical CDF percentiles.

For important variables such as \texttt{\_AGE80}, \texttt{\_BMI5}, and
\texttt{PHYSHLTH}, we also construct piecewise-linear spline features
$(x-\kappa_k)_+$ at empirical knots
$\kappa_k\in\{q_{0.05},\dots,q_{0.95}\}$. Finally, pairwise interactions are
generated among 13 major clinical variables, including \texttt{\_AGE80},
\texttt{SEX}, \texttt{GENHLTH}, \texttt{BPHIGH4}, \texttt{TOLDHI2},
\texttt{\_SMOKER3}, and \texttt{DIABETE3}.

% ============================================================
% METHODOLOGY
% ============================================================
\section{Methodology}

\subsection{Course Baselines}

We implemented all six required methods in \texttt{implementations.py}: MSE
regression using gradient descent and stochastic gradient descent, ordinary
least squares, ridge regression, logistic regression, and regularized logistic
regression.

Least squares is poorly matched to this task because its predictions are
unbounded and its squared-error objective does not directly model binary class
probabilities. Logistic regression instead models
$P(y=1\mid x)=\sigma(x^\top w+b)$ and forms the basis of our final pipeline.

\subsection{Regularized Logistic Optimization}

Our penalized objective is

\begin{equation}
\begin{aligned}
\min_{w,b}\quad
&\frac{1}{N}\sum_{i=1}^{N}
\Big[
\log(1+e^{x_i^\top w+b}) \\
&\qquad - y_i(x_i^\top w+b)
\Big]
+\lambda\|w\|_2^2+\mu\|w\|_1 .
\end{aligned}
\label{eq:objective}
\end{equation}

The intercept $b$ is not regularized. For smooth $L_2$ logistic regression
($\mu=0$), we implement L-BFGS~\cite{nocedal2006} with $m=12$ correction pairs and Armijo
backtracking. Step sizes are chosen by the line search rather than a fixed learning-rate schedule.
For Elastic Net ($\mu>0$), we use FISTA~\cite{beck2009fista} with adaptive momentum restarts and
the soft-thresholding operator

\begin{equation}
\operatorname{prox}_{\gamma\mu}(v)_j
=
\operatorname{sign}(v_j)
\max\!\left(|v_j|-\gamma\mu,0\right).
\end{equation}

Numerical stability is ensured through \texttt{logaddexp(0.0,z)} and stable
sigmoid computations.

\subsection{Threshold Optimization and Blending}

Because of class imbalance, the default threshold $\tau=0.5$ produces too many
false negatives. We therefore sweep all sorted out-of-fold probabilities and
use cumulative counts to find the threshold maximizing $F_1$ in
$O(N\log N)$ time.

The planned ensemble search compares probability averaging with averaging
in log-odds space:

\begin{equation}
\begin{aligned}
\operatorname{logit}(p_{\mathrm{blend}})
&=\sum_{m=1}^{M}\alpha_m\operatorname{logit}(p_m),\\
\alpha_m&\geq0,\qquad
\sum_{m=1}^{M}\alpha_m=1 .
\end{aligned}
\label{eq:blend}
\end{equation}

% ============================================================
% RESULTS
% ============================================================
\section{Experimental Results}

\subsection{Interim Validation Result}

The current run uses a stratified 80/20 split of the training data, with
262,508 fitting rows and 65,627 validation rows. The logistic-only search
reached $F_1=0.449939$ on this validation split. Candidate representations,
regularization settings, and decision thresholds are compared using these
validation labels, so this value is an interim tuning score rather than an
independent estimate of the final selected pipeline's performance.

\begin{table}[t]
\centering
\caption{Confirmed Interim Logistic-Only Result}
\label{tab:interim}
\begin{tabular}{lr}
\toprule
\textbf{Quantity} & \textbf{Value} \\
\midrule
Fitting rows & 262,508 \\
Validation rows & 65,627 \\
Validation $F_1$ & 0.449939 \\
\bottomrule
\end{tabular}
\end{table}

\subsection{Remaining Evaluation}

The run stopped before final cross-validation was completed. The remaining
steps are to finish the logistic-only search, generate three-fold out-of-fold
predictions, and select ensemble weights and a decision threshold. Those
out-of-fold scores will also be tuning scores if they guide selection.
The selected component models can then be refitted on all training rows to
produce a submission.

No final cross-validation score, ensemble composition, coefficient-sparsity
measurement, or final decision threshold is established by the interim run.
These quantities must be reported after the unfinished runs have completed.

% ============================================================
% CONCLUSION
% ============================================================
\section{Discussion \& Conclusion}

The current logistic-only pipeline combines survey preprocessing, nonlinear
feature representations, and regularization while retaining a linear model in
the expanded feature space. Its confirmed result is an interim validation
$F_1$ of $0.449939$. This supports continuing the regression-based search, but
does not establish which encoding, regularization setting, or blend will be
best after the remaining evaluation.

Final conclusions will require completion of the planned cross-validation
and ensemble-selection runs. The report will then be updated with the final
logistic-only results and a clear distinction between tuning scores and any
independent performance estimates.

\balance

\bibliographystyle{IEEEtran}
\bibliography{literature}

\end{document}
